\section{Geometry and degenerate dissipation for four-wave collisions with a cube cutoff}\label{geom:section}

We now turn to quadratic dispersion in three dimensions. The full four-leg collision
difference again connects conservation with dissipation: after grouping the momenta
into incoming and outgoing pairs, the entire dissipation becomes a pair-sum variance
over the admissible spherical directions. The pair sums at every center and radius
come from the same function and therefore jointly determine the five-dimensional
null space. Fix $R>0$ and let $D=[-R,R]^3$ be the momentum domain; all four legs
of every collision are restricted to this same cube. At its corners, the set of
admissible collisions degenerates and the collision frequency vanishes. We first
compute this rate of degeneracy, then use the frequency as a weight to establish
coercivity and an inverse for the linear equation. Finally, we prove the corner-layer
positivity and spatial product bounds needed in the nonlinear energy estimates.
To distinguish momentum components from collision-leg indices, write
$k=(k^{(1)},k^{(2)},k^{(3)})$ and set
\[
 \Psi(k)=(1,k^{(1)},k^{(2)},k^{(3)},|k|^2)^T,\qquad
 \Theta_+=\{\theta\in\R^5:\min_{k\in D}\theta\cdot\Psi(k)>0\},
 \qquad N_\theta=(\theta\cdot\Psi)^{-1}.
\]
Unless stated otherwise, the parameters lie in a fixed compact set
$\Theta\Subset\Theta_+$. Thus $N_\theta$ and its parameter and momentum
derivatives of any fixed order are uniformly bounded, and
$0<N_-\le N_\theta\le N_+<\infty$. Constants may depend on $R,\Theta$,
but not on the strong-collision coefficient introduced later.

\subsection{The full resonance measure and two integral kernels}

Write $\chi=\mathbf1_D$. For a momentum quadruple $(k_0,k_1,k_2,k_3)$,
indices $0,1$ denote the incoming legs and $2,3$ the outgoing legs. We abbreviate
$\phi_i=\phi(k_i)$ and $f_i=f(k_i)$, and define
\begin{align}
 \dd\Xi_D={}&\delta^3(k_0+k_1-k_2-k_3)
 \delta(|k_0|^2+|k_1|^2-|k_2|^2-|k_3|^2)
 \prod_{i=0}^3\chi(k_i)\,\dd k_i,\label{geom:measure}\\
 \Delta\phi={}&\phi_0+\phi_1-\phi_2-\phi_3.
\end{align}
The $\delta$ distributions here are defined by the coarea formulas below.
The four indicator functions jointly restrict a collision and remain present
after variables are eliminated. If $f$ also depends on time and space, all four
values $f_i=f(t,X,k_i)$ are taken at the same $(t,X)$.
Fix $k_0=k$ and set $k_2=k+x,k_3=k+y,k_1=k+x+y$. Then
\begin{equation}\label{geom:gamma}
\dd\Gamma_k=\frac12\delta(x\cdot y)
\chi(k)\chi(k+x)\chi(k+y)\chi(k+x+y)\,\dd x\dd y.
\end{equation}
The momentum constraint is automatic in these coordinates, and the energy
difference is $2x\cdot y$, which gives the factor $1/2$. A resonant quadruple
is therefore a rectangle with vertex $k$ and orthogonal sides $x,y$.
Relative to the output $k_0$, we call the diagonally opposite vertex $k_1$
the opposite leg, and the neighboring vertices $k_2,k_3$ the adjacent legs.
For $x\ne0$, the measure is
$\frac1{2|x|}\dd x\,\dd\mathcal H^2(y)|_{x^\perp}$, multiplied by all
the indicators above. No additional atom is assigned to $x=0$.
The admissible $x,y$ are bounded, and $|x|^{-1}$ is locally integrable
in three dimensions. This defines a finite positive measure, with
$\dd\Xi_D=\dd k\,\dd\Gamma_k$.
The collision operator is
\begin{equation}\label{geom:collision}
 \mathcal C(f)_0=\int
 (f_1f_2f_3+f_0f_2f_3-f_0f_1f_3-f_0f_1f_2)\,\dd\Gamma_{k_0}.
\end{equation}

\begin{lemma}[Boundedness, continuity, and conservation]\label{geom:basic}
Each trilinear term in \eqref{geom:collision} maps $C(D)^3$ continuously
into $C(D)$, with norm bounded by $C_R$; the same norm bound holds
on $L^\infty(D)$. All five moments satisfy
$\int_D\Psi\mathcal C(f)\,\dd k=0$.
If $N=N_\theta$, then $\mathcal C(N)=0$.
\end{lemma}
\begin{proof}
The coarea formula bounds the absolute value of each term by
$C_R\prod_j\|f_j\|_\infty\int_{|x|\le C_R}|x|^{-1}\dd x$.
To prove continuity of the output on the closed domain, fix $k_n,k\in D$
with $k_n\to k$. The output indicator is then identically $1$.
For almost every $x$, the point $k+x$ lies on no face of the cube,
and $x$ is parallel to no coordinate axis. On the plane $x^\perp$,
each face equation for $k+y$ or $k+x+y$ is a nonconstant affine equation,
so its boundary has zero area. Hence all indicators and continuous factors
converge pointwise. Dominated convergence, first on a fixed bounded planar
disc and then against $|x|^{-1}\dd x$, proves continuity.
The same argument applies when the output lies on the boundary; the excluded
set of $x$ has zero three-dimensional measure.

For positive $f$, the cubic polynomial equals
$f_0f_1f_2f_3\Delta(1/f)$. The full measure is invariant under exchanges
within either pair and under exchange of the two pairs. Thus, for every
bounded $\phi$,
\[
 \int\phi\mathcal C(f)=\frac14\int\prod_i f_i\,
                  \Delta(1/f)\Delta\phi\,\dd\Xi_D.
\]
Taking $\phi=\Psi_j$ gives conservation. This is a polynomial identity,
so it also holds for bounded real $f$: one may first add a positive constant
to $f$ and then compare coefficients. Finally,
$\Delta(N^{-1})=\theta\cdot\Delta\Psi=0$.
\end{proof}

Take $N=N_\theta$. Writing a perturbation as $f=N+Nu$, define the collision
frequency and the dissipative operator acting on the relative perturbation $u$ by
\begin{equation}\label{geom:linear}
 \nu_N(k)=N(k)^{-1}\int N_1N_2N_3\,\dd\Gamma_k,
 \qquad L_Nu=-N^{-1}D\mathcal C(N)[Nu].
\end{equation}
Its quadratic form is
\begin{equation}\label{geom:form}
 \mathfrak a_N(u,v)=\frac14\int\prod_iN_i\,
       \Delta(u/N)\overline{\Delta(v/N)}\,\dd\Xi_D.
\end{equation}
On unweighted $L^2(D)$, this is the form of a bounded nonnegative
self-adjoint operator, and
\begin{equation}\label{geom:kdecomp}
 L_N=\nu_N+K_{1,N}-2K_{2,N}.
\end{equation}
The two kernels correspond to integration over the opposite and adjacent legs.
Below, $k,p\in D$, and $\chi N$ denotes the extension of $N$ by zero outside $D$.
Fix $k_0=k,k_1=p$ and put $m=(k+p)/2,r=|k-p|/2$. The energy difference is
$2(r^2-|k_2-m|^2)$, whose normal derivative on the sphere has magnitude
$4r$. Hence
\begin{equation}\label{geom:k1}
 K_{1,N}(k,p)=\frac{|k-p|}{8}\int_{\mathbb S^2}
       (\chi N)(m+r\omega)(\chi N)(m-r\omega)\,\dd\omega.
\end{equation}
Fix instead $k_0=k,k_2=p$, and set $x=p-k,k_3=k+z,k_1=p+z$.
The energy difference is $2x\cdot z$, so
\begin{equation}\label{geom:k2}
 K_{2,N}(k,p)=\frac1{2|p-k|}\int_{z\perp(p-k)}
       (\chi N)(p+z)(\chi N)(k+z)\,\dd\mathcal H^2(z).
\end{equation}
Both kernels are nonnegative and symmetric, and
\begin{equation}\label{geom:pointkernel}
 K_{1,N}\le C|k-p|,\qquad K_{2,N}\le C|k-p|^{-1}.
\end{equation}
Their diagonal values do not affect the integral operators. The two adjacent
inputs give the same kernel $K_{2,N}$, which accounts for the factor $2$
in \eqref{geom:kdecomp}. Boundedness follows from the Schur bounds supplied
by \eqref{geom:pointkernel}; expanding the full squared difference gives
\eqref{geom:form}.

Pair coordinates give another representation of the dissipation.
For fixed center $V$ and radius $r$, the pair $V\pm r\sigma$ has total
momentum $2V$ and total energy $2(|V|^2+r^2)$; only the direction $\sigma$
varies. A collision therefore compares two directions at the same center
and radius, while the sharp cube cutoff determines which directions are admissible.

\begin{proposition}[Variance representation of the full pair sums]\label{geom:pairedvariance}
For $V\in\R^3$, $r>0$, and $\sigma\in\mathbb S^2$, set
\[
 A_{V,r}(\sigma)=(\chi N)(V+r\sigma)(\chi N)(V-r\sigma),
 \qquad Z_{V,r}=\int_{\mathbb S^2}A_{V,r}(\sigma)\dd\sigma.
\]
When $A_{V,r}(\sigma)>0$, define
\[
 H_u(V,r,\sigma)=\frac{u(V+r\sigma)}{N(V+r\sigma)}
                         +\frac{u(V-r\sigma)}{N(V-r\sigma)}.
\]
Set its value to zero in all other directions. If $Z_{V,r}>0$, write
$\overline H_u=Z_{V,r}^{-1}\int A_{V,r}H_u\dd\sigma$;
if $Z_{V,r}=0$, set $\overline H_u=0$.
For bounded $u,v$,
\begin{equation}\label{geom:varianceidentity}
 \mathfrak a_N(u,v)=
 \int_{\R^3}\int_0^\infty r^3Z_{V,r}
       \int_{\mathbb S^2}A_{V,r}(\sigma)
       (H_u-\overline H_u)\overline{(H_v-\overline H_v)}
                    \dd\sigma\dd r\dd V.
\end{equation}
The same identity extends continuously to the space $H_\nu$ defined below.
In particular, vanishing of the form requires that, for almost every fixed
$(V,r)$, the pair sum arising from the same $u$ be constant over all
admissible directions.
\end{proposition}
\begin{proof}
First eliminate the momentum constraint and write
$k_0=V+p,k_1=V-p,k_2=V+q$, so that $k_3=V-q$.
The absolute Jacobian determinant of this nine-dimensional transformation
is $8$, and the energy constraint is $2(|p|^2-|q|^2)$.
Passing to polar coordinates $p=r\sigma,q=s\tau$ and using the energy
constraint to eliminate $s$ gives
\[
 \dd\Xi_D=2r^3\prod_{i=0}^3\chi(k_i)\,
                      \dd V\dd r\dd\sigma\dd\tau.
\]
The full difference in \eqref{geom:form} is then
$H_u(\sigma)-H_u(\tau)$, and the two angular weights are
$A_{V,r}(\sigma)A_{V,r}(\tau)$. Expanding the product of the two differences
gives the cross term $(\int AH_u)\overline{(\int AH_v)}$. Thus
\[
 \frac12\int\!\!\int A(\sigma)A(\tau)
       [H_u(\sigma)-H_u(\tau)]
       \overline{[H_v(\sigma)-H_v(\tau)]}\dd\sigma\dd\tau
   =Z\int A(H_u-\overline H_u)
                 \overline{(H_v-\overline H_v)}\dd\sigma.
\]
This proves the coefficient and all cutoff conditions in
\eqref{geom:varianceidentity}. If $Z=0$, both sides vanish.
The single-leg marginal gives
$\int\prod_iN_i|u_i/N_i|^2\dd\Xi_D=\int\nu_N|u|^2$.
Hence both the uncentered pair-sum map and the map obtained by subtracting
its weighted mean are bounded on $H_\nu$. Approximation by bounded functions
gives the extension; all pair sums and their means are understood as
almost-everywhere representatives for the corresponding measures.
\end{proof}

\subsection{The precise degeneracy of the corner frequency}

Let $Z$ be the set of eight corners and
$r(k)=\min_{\xi\in Z}|k-\xi|_\infty$. For $\eta>0$, write
\[
 E_\eta(\xi)=\{k\in D:|k-\xi|_\infty\le\eta\},\qquad
 E_\eta=\bigcup_{\xi\in Z}E_\eta(\xi)=\{k\in D:r(k)\le\eta\}.
\]
Near a corner, use the three coordinate products $x^{(j)}y^{(j)}$
as variables. Their positive and negative densities can be computed
explicitly. The constraint that their sum vanish expresses the frequency
as convolutions for six sign patterns. The squared logarithm below comes
from two of the positive product densities.

\begin{proposition}[The squared-logarithmic corner frequency]\label{geom:frequency}
There exist $r_0,c_0,C_0>0$ such that, whenever $0<r(k)<r_0$,
\begin{equation}\label{geom:frequencyestimate}
 c_0r(k)^2\log^2\frac{2R}{r(k)}
 \le\nu_N(k)\le C_0r(k)^2\log^2\frac{2R}{r(k)}.
\end{equation}
The function $\nu_N$ is continuous on the closed cube, and its zero set
is exactly $Z$. In particular, for $p\ge0$,
$\nu_N^{-p}\in L^1(D)$ if and only if $p\le3/2$.
These conclusions are uniform on $\Theta$.
\end{proposition}
\begin{proof}
It suffices to compute the unweighted frequency $M(k)=\int\dd\Gamma_k$,
since $\nu_N\asymp M$. At the corner $(R,R,R)$, set
$L=2R,d_j=(R-k^{(j)})/L,e=\max_jd_j$, and put $x=L X,y=L Y$.
The full admissible domain in each coordinate is precisely
\[
 \mathcal P_d=\{X,Y\in[-1+d,d]:X+Y\in[-1+d,d]\}.
\]
Push forward the area measure of this planar domain under the product
$t=XY$. Denote its density on the positive half-line by $P_d(t)$,
and write its density on the negative half-line as $N_d(s)$, with $t=-s$.
They are, respectively,
\begin{align}
 P_d(t)&=T_{1-d}(t)+T_d(t),\quad
 T_b(t)=\log\frac{(b+\sqrt{b^2-4t})^2}{4t}\,
                  \mathbf1_{0<t<b^2/4},\label{geom:productpositive}\\
 N_d(s)&=2\log\frac{d(1-d)}s\,
                  \mathbf1_{0<s<d(1-d)},\quad t=-s.\label{geom:productnegative}
\end{align}
The formulas on the right are used only on the intervals specified by
the indicators and are defined to be zero elsewhere; also $T_0=N_0=0$.
The first formula comes from the two triangles on which the coordinates
have the same sign: solve $uv=t,u+v\le b$ and integrate $\dd u/u$.
The mixed-sign quadrant is exactly
$0\le X\le d,0\le-Y\le1-d$; its endpoints already satisfy the third
indicator, giving the second formula. There are no atoms on the axes.

For the positive densities extended by zero, define
$(P_{d_i}*P_{d_l})(s)=\int_0^sP_{d_i}(t)P_{d_l}(s-t)\dd t$.
When the three products sum to zero, their signs, outside the null
coordinate axes, must be one negative and two positive, or two negative
and one positive. Writing $d=(d_1,d_2,d_3)$, we have $M=L^4m(d)$,
and the six sign patterns give
\begin{align}
 m(d)={}&\frac12\sum_{j=1}^3\int_0^\infty
           N_{d_j}(s)(P_{d_i}*P_{d_l})(s)\,\dd s\notag\\
 &+\frac12\sum_{j=1}^3\int_0^\infty\!\int_0^\infty
       N_{d_i}(u)N_{d_l}(v)P_{d_j}(u+v)\,\dd u\dd v,
       \qquad\{i,l,j\}=\{1,2,3\}.\label{geom:sixsigns}
\end{align}
This identity can also be obtained directly in each orthant by the change
of variables $t_j=X_jY_j$. The absolute integrability bounds below
also determine the value of the convolution at zero energy.
For $e\le1/64$,
$P_{d_j}(t)\le2\log(1/t)$ and
$N_{d_j}(s)\le2\log(e/s)\mathbf1_{s<e}$.
Setting $t=sv$ gives
\[
 (P_{d_i}*P_{d_l})(s)
 \le Cs\int_0^1[\log(1/s)+|\log v|]
                    [\log(1/s)+|\log(1-v)|]\dd v
 \le Cs[1+\log(1/s)]^2.
\]
The first line of \eqref{geom:sixsigns} is therefore at most
$Ce^2\log^2(1/e)$. In the second line, set $u=eu',v=ev'$ and use
$\log_+(1/(u'+v'))\le\min(|\log u'|,|\log v'|)$
to obtain the bound $Ce^2[1+\log(1/e)]$.
All endpoint logarithms here are integrable.

For the lower bound, choose $d_j=e$ and restrict to
$s\in[e/8,e/4]$ and $t\in[s/3,2s/3]$. Then $N_e(s)\ge c$,
while the triangle side lengths in the other two coordinates satisfy
$1-d_i,1-d_l\ge63/64$. Hence
$T_{1-d_i}(t),T_{1-d_l}(s-t)\ge c\log(1/e)$.
This positive contribution gives $ce^2\log^2(1/e)$; the boundary
distances in the other two coordinates may vanish. Coordinate reflections
handle the remaining corners.

Continuity follows from the same coarea dominated-convergence argument
as in Lemma \ref{geom:basic}. At a noncorner point, at least one coordinate
$j$ lies in the open interval. For small $s>0$, set
$x^{(i)}=y^{(i)}=-\operatorname{sgn}(k^{(i)})s$ for $i\ne j$
(choose either sign when $k^{(i)}=0$), and set
$x^{(j)}=\sqrt2s,y^{(j)}=-\sqrt2s$.
All three moving legs lie strictly inside the cube, and
$x\cdot y=0,x\ne0$. A neighborhood on their resonance surface has positive
measure, so $M(k)>0$. At a corner, every coordinate product is nonnegative;
only the null axes can make their sum zero, and the preceding density
formula gives $M=0$. Compact sets away from the corners consequently
have a uniform positive lower bound. Finally, a corner shell has volume
$3r^2\dd r$, so integrability of the inverse frequency is equivalent to that of
\[
 \int_0^{r_0}r^{2-2p}\log^{-2p}(2R/r)\dd r.
\]
At the critical exponent $p=3/2$, this is
$\int\dd r/[r\log^3(2R/r)]$, which still converges; it diverges for
larger $p$.
\end{proof}

From now on, fix $n_*=(1+|k|^2)^{-1}$ and $\nu_* =\nu_{n_*}$, and write
$H_\nu=L^2(D,\nu_*\dd k)$. On compact parameter families,
$\nu_N\asymp\nu_*$. This weight is dictated by the fact that each
single-leg squared integral in the four-leg collision difference is
exactly $\int_D\nu_N|u|^2\dd k$. Dissipation therefore naturally controls
the $H_\nu$ norm, whereas the vanishing frequency permits unweighted
$L^2$ mass to concentrate at the corners; see Corollary \ref{geom:nogap}.

\subsection{The five collision invariants and weighted coercivity}

\begin{lemma}[The five-dimensional space of collision invariants]\label{geom:invariants}
If $\phi\in L^2_{\mathrm{loc}}(D^\circ)$ and
$\Delta\phi=0$ $\Xi_D$-almost everywhere, then
$\phi(k)=a+b\cdot k+c|k|^2$ almost everywhere.
\end{lemma}
\begin{proof}
The pair representation in Proposition \ref{geom:pairedvariance} shows
that a vanishing difference requires the pair sum to be constant over
admissible directions for almost every fixed $(V,r)$. These pair sums,
for different $(V,r)$, all come from the same function $\phi$.
In any strictly smaller concentric cube, translate all four legs together
and convolve in the translation parameter. The momentum and quadratic-energy
constraints are preserved under a common translation. Thus the smoothed function
$\phi_\delta$ still satisfies the collision relation, provided the four
test legs remain in a slightly smaller cube.
This operation is legitimate for the almost-everywhere relation with
respect to the resonance measure: write the legs as
$V+r\sigma,V-r\sigma,V+r\tau,V-r\tau$. Their measure is
$2r^3\dd V\dd r\dd\sigma\dd\tau$, whose density is strictly positive
for $r>0$. Fubini gives the almost-everywhere relation after common-translation
convolution, and smoothness then gives it for every admissible
$V,r,\sigma,\tau$. Hence
\[
 \phi_\delta(V+r\sigma)+\phi_\delta(V-r\sigma)
 =\phi_\delta(V+r\tau)+\phi_\delta(V-r\tau).
\]
Fix any interior points $x,y$ and write
$V=(x+y)/2$, $v=(x-y)/2$. Rotate $v$ in any coordinate plane, keeping
$|v|$ fixed, and differentiate the pair identity at zero rotation. This yields
\[
 [\partial_i\phi_\delta(x)-\partial_i\phi_\delta(y)](x^{(j)}-y^{(j)})
 =[\partial_j\phi_\delta(x)-\partial_j\phi_\delta(y)](x^{(i)}-y^{(i)}).
\]
These identities constrain the same gradient field simultaneously.
Choose the origin and three nearby points on the coordinate axes.
First compare increments between the axis points, then compare any $x$
with the origin and each axis point. This gives a single constant $\lambda$
such that $\nabla\phi_\delta(x)=b+\lambda x$.
Integration over the convex smaller cube gives the stated five-dimensional
family of quadratic polynomials; only first-order differentiability is used here.

Finally, take normalized smooth convolution kernels whose support radii
tend to zero with a fixed ratio of inner to outer radii.
The Lebesgue differentiation theorem gives pointwise convergence on a
single set of full measure. Equip the function space on this set with
the topology of pointwise convergence. The restrictions of the
five-dimensional polynomial family still form a finite-dimensional subspace,
which is closed in this Hausdorff topological vector space.
The limit on every smaller cube therefore remains in that family.
If two polynomials agree almost everywhere on a nonempty open cube,
continuity makes them equal everywhere there, and all their coefficients
agree. Exhaust $D^\circ$ by countably many smaller cubes and use the
Lebesgue-null boundary to obtain one set of coefficients on the full domain.
The same smoothing and limiting argument applies to locally integrable functions.
\end{proof}

\begin{lemma}[Weighted compactness]\label{geom:compact}
Define the weighted kernels
\[
 \widetilde K_{i,N}(k,p)
 =\nu_N(k)^{-1/2}K_{i,N}(k,p)\nu_N(p)^{-1/2},\qquad i=1,2.
\]
Both are Hilbert--Schmidt kernels, uniformly on $\Theta$.
\end{lemma}
\begin{proof}
By \eqref{geom:pointkernel}, it suffices to prove
$\iint\nu_*^{-1}(k)\nu_*^{-1}(p)|k-p|^{-2}\dd k\dd p<\infty$.
Choose $r_0>0$ sufficiently small as in Proposition \ref{geom:frequency},
and set $r_j=r_0 2^{-j}$. For $j\ge1$, combine the shells of equal depth
at all eight corners:
\[
 A_j=\{k\in D:r_{j+1}<r(k)\le r_j\},\qquad
 A_0=\{k\in D:r(k)>r_1\}.
\]
These sets partition the cube apart from its corners, and the closure
of $A_0$ stays away from the corners. Write
$a_j=\log^{-2}(2R/r_j)$. Proposition \ref{geom:frequency} and the
corner-layer volume give
\[
 |A_j|\le Cr_j^3,\qquad
 \sup_{A_j}\nu_*^{-1}\le C a_jr_j^{-2},\qquad
 a_j\asymp(1+j)^{-2}.
\]
After enlarging the constants, the zeroth layer satisfies the same bounds.
For any measurable set with $|A|\le Cr^3$, split the integral into
$|k-p|<r$ and its complement. Three-dimensional polar coordinates give directly
\[
 \sup_k\int_A|k-p|^{-2}\dd p
 \le |\mathbb S^2|r+r^{-2}|A|\le C'r.
\]
Thus, when $r_l\le r_j$, integration first over $A_j$ and then over $A_l$ gives
\[
 \int_{A_l}\int_{A_j}
 \frac{\dd k\dd p}{\nu_*(k)\nu_*(p)|k-p|^2}
 \le C a_ja_l\frac{r_l}{r_j}\le C a_ja_l.
\]
Treat the reverse order symmetrically and use $\sum_j a_j<\infty$.
This controls every pair of layers, including neighboring layers,
layers at different corners, and the interior region.
The corners form a null set and do not affect the integrals.
\end{proof}

The null-space classification and weighted compactness together yield
a positive lower bound on the complement of the five-dimensional null space.
Two notions of orthogonality must be distinguished: dissipation first controls
the weighted distance to the null space, whereas the microscopic constraint
uses the unweighted $L^2$ projection. We prove both statements and their relation.

\begin{proposition}[Coercivity with a degenerate weight]\label{geom:coercivity}
The form \eqref{geom:form} is defined for every $u\in H_\nu$, and
\begin{equation}\label{geom:weightedcoercivity}
 \mathfrak a_N(u,u)\ge c_\Theta
       \inf_{\psi\in\Span\{N\Psi\}}\|u-\psi\|_{H_\nu}^2.
\end{equation}
Let $P_N$ be the orthogonal projection in unweighted $L^2(D)$ onto
$\Span\{N\Psi\}$, and let $Q_N=I-P_N$. This projection extends to $H_\nu$, and
\begin{equation}\label{geom:microcoercivity}
 P_Nu=0\quad\Longrightarrow\quad
 c_\Theta\|u\|_{H_\nu}^2\le\mathfrak a_N(u,u)
                       \le C_\Theta\|u\|_{H_\nu}^2.
\end{equation}
\end{proposition}
\begin{proof}
Four-leg symmetry gives each single-leg integral as
$\int\prod N_i |u_i/N_i|^2\dd\Xi_D=\int\nu_N|u|^2$.
Every single-leg marginal is absolutely continuous with respect to volume.
Thus the volume almost-everywhere class of $u$ determines its value on
each leg, without assigning additional traces on boundaries or diagonals.
The form is therefore defined on all of $H_\nu$ and is bounded above by
$4\int\nu_N|u|^2$. Put $v=\sqrt{\nu_N}u$.
The full difference defines a bounded operator
$\mathcal Rv=\frac12\Delta(\nu_N^{-1/2}v/N)$ with values in
$L^2(\prod_iN_i\dd\Xi_D)$, whose squared norm is the form.
First expand for bounded $u$ supported away from the corners, then pass
to the limit in $H_\nu$. Lemma \ref{geom:compact} gives
\[
 \mathcal R^*\mathcal R=I+\widetilde K_{1,N}-2\widetilde K_{2,N}.
\]
By Lemma \ref{geom:invariants}, a zero form implies
$u\in\Span\{N\Psi\}$. Indeed, $u\in H_\nu$ belongs to $L^2$ on
every smaller cube, which suffices for that lemma; no global unweighted
$L^2$ assumption is made.
If there were no positive lower bound on the orthogonal complement of
$\sqrt{\nu_N}\Span\{N\Psi\}$, choose unit vectors $v_j$ whose form tends
to zero. The bounded nonnegative operator $T=\mathcal R^*\mathcal R$
satisfies $\|Tv_j\|^2\le\|T\|\langle Tv_j,v_j\rangle\to0$.
Compactness of $T-I$ gives a strongly convergent subsequence of
$v_j=Tv_j-(T-I)v_j$. Its limit is a nonzero vector both in the null space
and orthogonal to it, a contradiction.

For uniformity in the parameter, put $z=n_*u/N$. Then
$\Delta(z/n_*)=\Delta(u/N)$, and comparison of the positive weights
bounds $\mathfrak a_N(u,u)$ below by a constant multiple of
$\mathfrak a_{n_*}(z,z)$. Multiplication by $N/n_*$ is uniformly bounded
in both directions and identifies the two null spaces exactly.
The constant in \eqref{geom:weightedcoercivity} is consequently uniform.
Write $\kappa_N=N\Psi$ for the column vector of the five natural zero modes.
With $e_j$ the $j$th standard basis vector of $\R^5$, choose the real
orthonormal basis
\[
 M_N=\int_D N^2\Psi\Psi^T\dd k,\qquad
 \kappa_j(k)=N(k)\Psi(k)^TM_N^{-1/2}e_j,\quad 1\le j\le5.
\]
The matrix $M_N$ is strictly positive definite, with a positive lower
bound for its smallest eigenvalue on compact parameter sets.
The basis therefore depends smoothly on the parameter, and it and its
parameter derivatives of every fixed order are uniformly bounded.
Since $\kappa_j/\sqrt{\nu_*}\in L^2$, the projection $P_N$ is uniformly
bounded on $H_\nu$. If $P_Nu=0$ and $\psi$ is the best weighted
approximation from the null space, then $\psi=-P_N(u-\psi)$.
Hence $\|u\|_{H_\nu}\le C\|u-\psi\|_{H_\nu}$, proving the final assertion.
\end{proof}

\begin{corollary}[Absence of an unweighted spectral gap]\label{geom:nogap}
The collision operator has no strictly positive unweighted spectral gap
on $Q_NL^2(D)$.
\end{corollary}
\begin{proof}
Take an $L^2$-normalized function $g_\eta$ supported in a radius-$\eta$
neighborhood of one corner. This neighborhood has volume $O(\eta^3)$,
and the null-space basis is bounded, so
$\|P_Ng_\eta\|_2\le C\eta^{3/2}$. Moreover,
\[
 \mathfrak a_N(Q_Ng_\eta,Q_Ng_\eta)
 =\mathfrak a_N(g_\eta,g_\eta)
 \le C\sup_{E_\eta}\nu_*\longrightarrow0,
 \qquad\|Q_Ng_\eta\|_2\longrightarrow1.
\]
This is the required microscopic unit sequence with vanishing dissipation.
\end{proof}

\subsection{Absolute row estimates and the continuous weighted inverse}

Write $K_N=K_{1,N}-2K_{2,N}$ and $K_{\mathrm a,N}=K_{1,N}+2K_{2,N}$.
\begin{proposition}[Row bounds, squared-row bounds, and vanishing at corners]\label{geom:rows}
Uniformly in the parameter,
\begin{align}
 \int K_{\mathrm a,N}(k,p)\dd p&\le C\nu_*(k),\label{geom:rowplain}\\
 \int K_{\mathrm a,N}(k,p)^2\dd p&\le C\sqrt{\nu_*(k)},\label{geom:rowsquare}\\
 \sup_k\int\frac{K_{\mathrm a,N}(k,p)}{\nu_*(p)}\dd p&\le C.\label{geom:rowweighted}
\end{align}
The weighted row $k\mapsto K_{i,N}(k,\cdot)/\nu_*(\cdot)$ is continuous
in $L^1(D)$, and
\begin{equation}\label{geom:rowcorner}
 \sup_{\theta\in\Theta,\,k\in E_\eta}
     \int\frac{K_{\mathrm a,N}(k,p)}{\nu_*(p)}\dd p
       =:\psi(\eta)\longrightarrow0.
\end{equation}
\end{proposition}
\begin{proof}
Each kernel satisfies $\int N(p)K_{i,N}(k,p)\dd p=N(k)\nu_N(k)$,
so \eqref{geom:rowplain} retains the frequency factor.
For the squared-row bound, \eqref{geom:pointkernel} gives $C\rho$
on $|k-p|\le\rho$; outside this set, one pointwise kernel bound
and \eqref{geom:rowplain} give $C\nu_*(k)(1+\rho^{-1})$.
Take $\rho\asymp\sqrt{\nu_*(k)}$, capped at a fixed radius.
At the corners, nonnegativity and zero frequency make the rows exactly zero.

Near a corner $\xi$, write $r=|p-\xi|$ and $L(r)=\log(C_R/r)$.
The key estimate is the uniform Newton-potential tail bound
\begin{equation}\label{geom:newtontail}
 \sup_{k\in D}\int_{|p-\xi|<\eta}
 \frac{1+|k-p|^{-1}}{\nu_*(p)}\dd p\le\frac C{L(\eta)}.
\end{equation}
Set $\rho=|k-\xi|$. If $\rho>2\eta$, use the lower bound on the distance
and $\int_{r<\eta}\nu_*^{-1}\le C\eta/L(\eta)^2$.
For $0<\rho\le2\eta$, split into $r<\rho/2$,
$\rho/2\le r\le2\rho$, and $2\rho<r<\eta$. The first two contributions
are each bounded by $C/L(\rho)^2$, and the last by
$C\int_{2\rho}^{\eta}\dd r/[rL(r)^2]$.
For $\rho=0$, integrate directly down to zero.
Away from the corners, the inverse frequency is bounded, proving
\eqref{geom:rowweighted}; the remaining contribution from a small diagonal
ball is $C\rho^2$.

The spherical formula for $K_1$ is continuous off the diagonal, and its
linear upper bound gives continuous vanishing on the diagonal.
For $K_2$, fix an output $k_0$ and exclude the three lines
$p\in k_0+\R e_j$ and the point $p=k_0$ from the input variable.
Near any remaining $p$, orthonormal coordinates for $(p-k)^\perp$
can be chosen continuously. The preimage of each cutoff face is either
a line in the plane or the empty set, and therefore has zero area.
Dominated convergence thus gives kernel continuity for each such fixed $p$.
The excluded input set has zero three-dimensional measure.
First remove corner inputs using \eqref{geom:newtontail}, then remove a
small diagonal ball, and finally use the same bounded majorant on the
remaining domain. This proves continuity of the entire row in $L^1$,
including outputs on faces, edges, and corners.
Pointwise continuity of kernel values is used only for inputs outside
the specified null set. Compactness of the closed domain and parameter
set gives uniform continuity. Since the corner rows themselves vanish,
\eqref{geom:rowcorner} follows.
\end{proof}

Use the Banach spaces
\[
 \mathcal X=L^\infty(D),\qquad
 \mathcal Y=\{u:\nu_*u\in L^\infty(D)\},
 \qquad\|u\|_{\mathcal Y}=\|\nu_*u\|_\infty.
\]
Functions are understood as almost-everywhere equivalence classes;
membership in $\mathcal Y$ does not presuppose unweighted $L^2$ integrability.
Since $\nu_*^{-1}\in L^1$,
$\mathcal Y\subset H_\nu\cap L^1(D)$, so $P_N$ remains defined
on these functions. The weight allows solutions to grow near the corners;
the uniformly controlled quantity in the equation is $\nu_*u$.
This is the space used for the cell problem $L_Nu=F$ below.

\begin{proposition}[The normalized weighted inverse]\label{geom:cell}
For $F\in Q_N\mathcal X$, there is a unique normalized solution
$u\in Q_N\mathcal Y$ of $L_Nu=F$, and
\begin{equation}\label{geom:cellbound}
 \|u\|_{\mathcal Y}+\mathfrak a_N(u,u)^{1/2}
   \le C\|F\|_\infty.
\end{equation}
If $F$ is continuous, $\nu_Nu$ has a continuous representative.
The following absolute row bound also holds:
\begin{equation}\label{geom:absoluteformrow}
 \mathop{\mathrm{ess\,sup}}_{k\in D}N(k)\int N_1N_2N_3
             |\Delta(u/N)|\,\dd\Gamma_k\le C\|F\|_\infty.
\end{equation}
Here $u$ is an almost-everywhere equivalence class in the weighted space.
The continuity assertion specifies only a representative of $\nu_Nu$,
not corner values for $u$ itself.
\end{proposition}
\begin{proof}
Absorb the diagonal multiplication into the unknown to write the
equation as the identity plus a compact operator. The five-dimensional
conservation null space corresponds to five compatibility conditions.
A finite-rank term removes that null space; after inversion, its
contribution vanishes for microscopic right-hand sides.
Set $v=\nu_Nu$ and $J_Nv=K_N(\nu_N^{-1}v)$.
Proposition \ref{geom:rows} and the Arzel\`a--Ascoli theorem show that
$J_N:\mathcal X\to C(D)$ is compact.
We first identify the integral operator with the weighted form.
For $u\in\mathcal Y$ and bounded $h$, the symmetric absolute kernel satisfies
\[
 \iint K_{\mathrm a,N}(k,p)|u(p)h(k)|\dd k\dd p
       \le C\|h\|_\infty\int\nu_*|u|\dd p<\infty.
\]
Truncate $u$ both near the corners and in amplitude. The resulting
$u_n$ converge to $u$ in $H_\nu$, and
$\int\nu_*|u_n-u|\to0$. Continuity of the form, the absolute integral
bound above, and dominated convergence for the diagonal term give
\[
 \mathfrak a_N(u,h)=\int_D(\nu_Nu+K_Nu)\overline h\dd k.
\]
This passage to the limit requires only convergence of the pairings,
not $\nu_Nu_n\to\nu_Nu$ in $L^\infty$.
If $(I+J_N)v=0$, then $u=v/\nu_N\in H_\nu\cap L^1$, and the identity
gives $\mathfrak a_N(u,h)=0$ for all bounded tests. Their density in
$H_\nu$ implies $u\in\Span\{N\Psi\}$.
Conversely, these null vectors clearly give zero. Thus
$\Ker(I+J_N)=\nu_N\Span\{N\Psi\}$ has dimension five.

Using the basis $\kappa_j$ defined above, set
\[
 \mathcal A_{N,0}v=(I+J_N)v+
           \sum_{j=1}^5\kappa_j\int\kappa_j\frac v{\nu_N}\dd k.
\]
The symmetric kernel satisfies $K_N\kappa_j=-\nu_N\kappa_j$, so
$\int\kappa_j(I+J_N)v=0$. The absolute row bounds justify these
exchanges of integration.
If $\mathcal A_{N,0}v=0$, first pair with each $\kappa_j$ to show that
all normalization integrals vanish, then use the five-dimensional
null-space description to obtain $v=0$.
This operator is the identity plus a compact operator, so the Fredholm
alternative gives a bounded inverse.
One can also reduce this directly to finite-dimensional linear algebra:
the compact image of the kernel operator is equicontinuous in $C(D)$,
and a finite partition of unity approximates it in operator norm by
finite-rank operators. Absorb the small remainder into an invertible
identity-plus-small operator. For the remaining identity-plus-finite-rank
operator, injectivity is equivalent to surjectivity.
The proof applies both to $C(D)$ and to $\mathcal X$, since null vectors
are automatically continuous; it does not use norm density of $C(D)$
in $L^\infty$.
For microscopic $F$, every normalization integral of the solution
vanishes automatically, and the solution satisfies the unmodified equation.
Its energy equals $\int Fu$, so $\nu_N^{-1}\in L^1$ gives
\eqref{geom:cellbound}. Estimating all four terms separately yields
\[
 N(k)\int N_1N_2N_3|\Delta(u/N)|\dd\Gamma_k
 \le N(k)\left(\nu_N(k)|u(k)|+
                         \int K_{\mathrm a,N}(k,p)|u(p)|\dd p\right),
\]
which proves \eqref{geom:absoluteformrow}.
Uniformity in the parameter will be established in the next proof.
\end{proof}

\begin{proposition}[A uniform resolvent bound]\label{geom:resolvent}
For $s\ge1$, the collision operator satisfies
\begin{equation}\label{geom:resolventbound}
 \|(I+sL_N)^{-1}F\|_{\mathcal X}
   +s\|(I+sL_N)^{-1}F\|_{\mathcal Y}\le C\|F\|_{\mathcal X},
 \qquad F\in Q_N\mathcal X.
\end{equation}
Continuous inputs give continuous outputs.
\end{proposition}
\begin{proof}
For $0\le z\le1$, define
\[
 \mathcal A_{N,z}v=v+K_N\frac v{\nu_N+z}
        +\sum_j\kappa_j\int\kappa_j\frac v{\nu_N+z}\dd k.
\]
Each is the identity plus a compact operator.
Operator-norm continuity at $z=0$ follows from \eqref{geom:newtontail}:
first remove corner inputs uniformly, after which
$z/(\nu_N+z)$ tends to zero uniformly on the remaining domain.
For the finite-rank part, use $\nu_N^{-1}\in L^1$.
Parameter derivatives of each fixed order satisfy
$|\partial_\theta^\alpha\nu_N|\le C_\alpha\nu_*$ and
$|\partial_\theta^\alpha K_N|\le C_\alpha K_{\mathrm a,n_*}$,
which also gives operator-norm continuity in the parameter.

For $z>0$, if $\mathcal A_{N,z}v=0$, pairing with $\kappa_j$ gives
$(1+z)\int\kappa_jv/(\nu_N+z)=0$.
Consequently $u=v/(\nu_N+z)$ satisfies $(z+L_N)u=0$.
It belongs to $L^2$, and nonnegativity forces $u=0$.
The null space at $z=0$ was already excluded in Proposition \ref{geom:cell}.
The inverse norms are therefore uniformly bounded on the compact set
$\Theta\times[0,1]$. This also proves uniformity in \eqref{geom:cellbound}.

For microscopic $F$, solve $\mathcal A_{N,z}v=F$.
The same pairings make the finite-rank term vanish. Thus
$(z+L_N)^{-1}F=v/(\nu_N+z)$, with unweighted norm at most
$Cz^{-1}\|F\|_\infty$ and weighted norm at most $C\|F\|_\infty$.
Set $z=1/s$ and multiply by $z$ to obtain the assertion.
Continuity follows from the same inverse construction on $C(D)$;
there is no need to treat the unbounded multiplier $\nu_N^{-1}$
as a bounded similarity transformation.
\end{proof}

Under parameter differentiation, the space $Q_N\mathcal X$ moves with $N$.
The following change of coordinates identifies all five-dimensional
null spaces with a fixed reference null space, so that subsequent
derivatives can be taken on fixed Banach spaces.

\begin{lemma}[Parameter coordinates with a fixed five-dimensional null space]\label{geom:frame}
There is a real unitary operator $U_N$ on $L^2$ satisfying
$U_NP_*=P_NU_N$, where $P_*=P_{n_*}$ and $Q_*=I-P_*$.
The operator $U_N-I$ has finite rank, and $U_N,U_N^*$ and their parameter
derivatives of every fixed order are uniformly bounded on
$\mathcal X,\mathcal Y,H_\nu$.
Set $\widetilde L_N=U_N^*L_NU_N$. Its null space is fixed, and
all parameter derivatives of fixed order map $\mathcal Y$ boundedly into
$\mathcal X$.
\end{lemma}
\begin{proof}
The overlap Gram matrix of the two null spaces is
$\int Nn_*\Psi\Psi^T$, which is strictly positive definite.
Their largest principal angle is therefore strictly less than $\pi/2$,
and on the compact parameter set $\|P_N-P_*\|\le r<1$.
Let $E=P_N-P_*$ and choose
\[
 U_N=\big[(I-P_N)(I-P_*)+P_NP_*\big](I-E^2)^{-1/2}.
\]
Projection algebra gives $T^*T=TT^*=I-E^2$ for the first factor $T$,
as well as $TP_*=P_NT$; moreover, $E^2$ commutes with both projections.
The displayed operator is therefore unitary and has the required
intertwining property.
For uniform bounds on the different spaces, express $P_N$ by its
five-dimensional Gram formula. The integral kernels of its parameter
derivatives are finite sums of tensor products of uniformly bounded functions.
Since $\nu_*^{-1}\in L^1$, for
$\mathcal E=\mathcal X,\mathcal Y,H_\nu$ and every fixed multi-index
$\alpha$, the maps
\[
 \partial_\theta^\alpha E:\mathcal E\longrightarrow L^2,
 \qquad
 \partial_\theta^\alpha E:L^2\longrightarrow\mathcal X
\]
are uniformly bounded. These estimates require no inversion of the
combined generating family $\{N\Psi,n_*\Psi\}$.
Set
\[
 h(z)=\frac{(1-z)^{-1/2}-1}{z},\qquad h(0)=\frac12.
\]
On $L^2$, $\|E^2\|\le r^2<1$, so the power series for $h(E^2)$
and all its parameter derivatives of fixed order converge uniformly
in operator norm. The identity
\[
 (I-E^2)^{-1/2}-I=E\,h(E^2)\,E
\]
together with the two endpoint mappings above shows that the nonidentity
part of the inverse square root and its parameter derivatives are
uniformly bounded from each $\mathcal E$ to $\mathcal X$.
Since $\mathcal X$ embeds continuously into all three spaces and $T-I$
has the same kind of finite-rank integral kernel, $U_N,U_N^*$ and all
the specified parameter derivatives have the asserted bounds.
Their nonidentity parts have finite rank, and the estimates remain
valid as $N$ approaches $n_*$.
Finally, $L_N$ and its parameter derivatives are bounded from
$\mathcal Y$ to $\mathcal X$: for the diagonal term use
$|\partial^\alpha\nu_N|\le C\nu_*$, and for the remaining terms use
\eqref{geom:rowweighted}. The Leibniz rule gives the assertion for the
conjugated operator.
\end{proof}

\begin{lemma}[Two-norm estimates for the cubic remainder]\label{geom:remainder}
Write
\[
 N^{-1}\mathcal C(N(1+q))=-L_Nq+\mathcal R_N(q).
\]
Then $P_N\mathcal R_N(q)=0$. If $\|q_i\|_{\mathcal X}\le\kappa\le1$,
\begin{align}
 \|\mathcal R_N(q)\|_{\mathcal X}&\le C\kappa\|q\|_{\mathcal Y},\notag\\
 \|\mathcal R_N(q_1)-\mathcal R_N(q_2)\|_{\mathcal X}
     &\le C\kappa\|q_1-q_2\|_{\mathcal Y}.\label{geom:remainderbound}
\end{align}
In each term of the second and third state derivatives and of parameter
derivatives of any fixed order, exactly one input may be placed in
$\mathcal Y$, with the others in $\mathcal X$, with uniform constants.
\end{lemma}
\begin{proof}
Each remainder term is a product of two or three function values on the
same quadruple. Place only one input in the weighted space $\mathcal Y$
and bound the others in the unweighted supremum norm.
The remaining integral is controlled precisely by the diagonal frequency
or by one of the two absolute row bounds.
Let $(s_0,s_1,s_2,s_3)=(1,1,-1,-1)$. Expand the four cubic monomials
term by term. The constant terms cancel because $\sum_i s_i/N_i=0$,
and the linear terms satisfy
\[
 \sum_i\frac{s_i}{N_i}\sum_{j\ne i}q_j=-\Delta(q/N).
\]
Thus the linear part is exactly $-L_Nq$.
All terms in the remainder can be written as
\[
 \mathcal R_N(q)_0=\int N_1N_2N_3
 \sum_{i=0}^3\frac{s_i}{N_i}
 \left(\sum_{\substack{j<l\\j,l\ne i}}q_jq_l
                         +\prod_{j\ne i}q_j\right)\dd\Gamma_{k_0}.
\]
There are twelve quadratic terms and four cubic terms, all retaining
the four-leg indicators.
In any term, designate one input $v_j$ and bound the other inputs in
the unweighted supremum norm. Since $N_j/N_i$ is uniformly bounded,
the remaining absolute integral is at most a constant multiple of
\[
 \begin{cases}
  \nu_N(k_0)|v_0(k_0)|,&j=0,\\
  \displaystyle\int K_{1,N}(k_0,p)|v_1(p)|\dd p,&j=1,\\
  \displaystyle\int K_{2,N}(k_0,p)|v_j(p)|\dd p,&j=2,3.
 \end{cases}
\]
This gives the bound from a $\mathcal Y$ input to a $\mathcal X$ output
for each position in the product.
Expand a difference by changing one factor at a time. Exactly one factor
is the difference of the two states; designate it as $v_j$.
At least one remaining input contributes a factor $\kappa$.
The two stated estimates follow from \eqref{geom:rowweighted}.
Parameter derivatives act only on the weights $N$ and the output
factor $N^{-1}$, not on the cutoff indicators, so the same proof applies
at every fixed order. Finally, five-moment conservation gives the
microscopic range. No product bound
$\mathcal Y\times\mathcal Y\to\mathcal X$ is used.
\end{proof}

\subsection{Linearization and spatial products in a thin corner layer}

We now allow the distribution to differ from $N$ by a finite amplitude
inside the corner layer, requiring closeness to $N$ only outside it.
To control the sign of the linearization, we first estimate, for a fixed
output, the conditional mass of the other legs entering a corner layer.
Opposite and adjacent legs play different geometric roles and must be
treated separately.
\begin{lemma}[Conditional corner-layer mass by leg position]\label{geom:roles}
There exist $\rho_0,C>0$, depending only on $R$, such that, for
$0<\eta\le\rho_0$ and all $k\in D$,
\begin{equation}\label{geom:oppositerole}
 \int\mathbf1_{k_1\in E_\eta}\dd\Gamma_k\le C\eta M(k).
\end{equation}
For each corner $\xi$, if $k\in E_{\rho_0}(\xi)$, then, for $i=2,3$,
\begin{equation}\label{geom:wrongcorner}
 \int\mathbf1_{k_i\in E_\eta\setminus E_\eta(\xi)}\dd\Gamma_k
       \le C\eta M(k).
\end{equation}
No small conditional-mass assertion is made when an adjacent leg and
the output lie in the same corner layer.
\end{lemma}
\begin{proof}
Choose $\rho_0$ sufficiently small that the corresponding dimensionless
corner depth lies in the range of Proposition \ref{geom:frequency},
and that an input at a different corner has one coordinate uniformly
separated from zero. This neighborhood is fixed before $\eta$ is chosen.
Use the scaling in \eqref{geom:productpositive}--\eqref{geom:sixsigns},
and denote the output corner depth by $d=\max d_j$.
On resonance, each $|X^{(j)}Y^{(j)}|\le2d$.
First observe that if the product density of any one coordinate is bounded
by $B$ in this range, while the other two densities are unchanged, the
sum of all six contributions is at most
$CBd^2[1+\log(1/d)]^2$.
If the designated coordinate is the sole negative one, use directly
$B\int_0^{2d}s[1+\log(1/s)]^2\dd s$.
If it is positive, first integrate the other positive coordinate to obtain
$Cs[1+\log(1/s)]$, then integrate against the negative density
$\log(d/s)$.
For the two-negative patterns, divide both negative variables by $d$;
the resulting endpoint terms contain only integrable logarithms.
This proves the uniform bounded-density replacement estimate.

If an adjacent input lies at a different corner, one coordinate $X$ or
$Y$ lies in a far-end strip of width $\eta$ and has absolute value
bounded below by a fixed positive constant. The product density of this
strip is bounded by $\int_{\mathrm{strip}}\dd X/|X|\le C\eta$;
the third indicator can only shrink the domain.
Divide the replacement estimate by
$M(k)\ge cd^2[1+\log(1/d)]^2$ to obtain \eqref{geom:wrongcorner}.

For the opposite leg, when the output has distance at least $\eta$
from every corner, the spherical kernel is uniformly bounded.
The input corner layer has volume $O(\eta^3)$ and
$M(k)\ge c\eta^2$, giving $C\eta$.
If the output corner depth satisfies $d\le\eta$, distinguish the input
corners. An opposite leg at the same corner forces both adjacent legs,
by momentum conservation and all the indicators, into a cube of side
$2\eta$ at that corner. Its linear scale is $2\eta$, and the dimensionless
output corner depth within it is $d/(2\eta)$.
The six-dimensional volume and the quadratic-energy $\delta$ distribution
together scale by $(2\eta)^4$.
Apply the unweighted version of Proposition \ref{geom:frequency},
absorbing the upper bound on compact sets away from corners into the
constant. Its upper bound then applies throughout
$0<d/(2\eta)\le1/2$ and bounds the frequency of this small cube by
$C\eta^2d^2[1+\log(\eta/d)]^2$.
Dividing by the full frequency gives $O(\eta^2)$.
If the opposite leg is at a different corner, one coordinate $X+Y$
lies in a far-end strip of width $\eta$.
In the quadrant where both coordinates are negative, set
$s=-X-Y,t=XY$. Then $s\ge c>0,|t|\le2d$, and the root separation
$\sqrt{s^2-4t}$ is uniformly nonzero.
In a mixed-sign quadrant, the root separation is $\sqrt{s^2+4|t|}$.
The product density is therefore at most $C\eta$, and the replacement
estimate applies again. At the corners themselves, every row vanishes,
so no division by zero frequency is needed.
\end{proof}

\begin{proposition}[Linearized coercivity under two-sided bounds in the corner layer]\label{geom:localized}
Fix $0<a\le1\le b<\infty$. There exist $\kappa_0,\eta_0,\gamma>0$
such that, if
\[
 a\le f/N\le b,\qquad |f/N-1|\le\kappa\quad\text{on }D\setminus E_\eta,
 \qquad\kappa\le\kappa_0,\quad\eta\le\eta_0,
\]
then, for every $h\in L^2(D)$ with $P_Nh=0$,
\begin{equation}\label{geom:localizedcoercivity}
 \RePart\left\langle h,-N^{-1}D\mathcal C(f)[Nh]\right\rangle
       \ge\gamma\|h\|_{H_\nu}^2.
\end{equation}
Moreover, the full operator can be written as $\nu_f+K_f$, with
\[
 c\nu_N\le\nu_f\le C\nu_N\quad(D\setminus Z),\qquad
 |K_f(k,p)|\le C K_{\mathrm a,N}(k,p),\qquad
 \int|K_f(k,p)|^2\dd p\le C\sqrt{\nu_*(k)}.
\]
This quadratic form also extends continuously to $H_\nu$, and the
inequality remains valid there. Its diagonal term is the integral of
$\nu_f|h|^2$, and its off-diagonal term is defined by the bounded form
of the weighted kernel. Equivalently, it is the limit obtained by
approximating with bounded functions in $H_\nu$; a function belonging
only to $H_\nu$ is not treated directly as both entries of an unweighted
$L^2$ inner product.
The corner rows remain zero.
All bounds also hold after averaging $D\mathcal C$ between positive
states satisfying the same conditions. The statement retains the full
gain and loss terms and does not require the right null space for
general $f$ to remain $N\Psi$.
\end{proposition}
\begin{proof}
There are three steps. Conditional-mass estimates first preserve a
positive lower bound for the diagonal frequency. The off-diagonal part
is then treated as a small perturbation of the compact equilibrium kernel.
Finally, weak convergence separates the mass of the limit from mass
that may concentrate at the corners.
The diagonal coefficient is
\[
 \nu_f(k_0)=\int a_0\dd\Gamma_{k_0},\qquad
 a_0=f_1f_3+f_1f_2-f_2f_3=f_1f_3+f_2(f_1-f_3).
\]
An individual $a_0$ may change sign. Set $r=\sqrt\eta$ and take $\eta$
sufficiently small that $4r\le\rho_0$, with $\rho_0$ from
Lemma \ref{geom:roles}.
When the output is at least distance $r$ from the corners, the spherical
and planar kernels, together with $M(k)\ge cr^2$, bound the conditional
mass of each corner-layer input by $C\eta^3/r^3$. Consequently
$\nu_f/\nu_N=1+O(\kappa+\eta^3/r^3)$.
When the output is within distance $r$ of $\xi$, mark the part where
$k_1\in E_{4r}$ or an adjacent input lies at another corner.
Lemma \ref{geom:roles} bounds its conditional mass by $O(r+\eta)$.
On the complement there is at most one bad nonoutput leg, and it can
only be an adjacent leg at the same corner: two such adjacent bad legs
would force the opposite leg into $E_{4r}$.
If the bad leg is $k_2$, the other two lie outside the layer, and
$|k_1-k_3|=|k_2-k_0|\le C(r+\eta)$. Thus
$|f_1-f_3|\le C(\kappa+r)$ and
$a_0\ge N_-^2-C_b(\kappa+r)$.
Exchange the roles if the bad leg is $k_3$; when there is no bad leg,
use the RJ identity. On the marked part, use only the fixed absolute
upper bound. This gives
\begin{equation}\label{geom:fulldiagonal}
 0<c\nu_N\le\nu_f\le C\nu_N
 \quad\text{on }D\setminus Z.
\end{equation}
Moreover, as $\kappa,\eta\to0$, $\nu_f/\nu_N\to1$ almost everywhere.

The other three coefficients are
\[
 a_1=f_0f_3+f_0f_2-f_2f_3,\quad
 a_2=f_1f_3+f_0f_3-f_0f_1,\quad
 a_3=f_1f_2+f_0f_2-f_0f_1,
\]
and enter the kernel with signs $+a_1,-a_2,-a_3$ and the corresponding
factors $N_i/N_0$.
The fixed positive two-sided bounds control their absolute values by
$K_{\mathrm a,N}$, and the squared-row bound follows from
Proposition \ref{geom:rows}.
Denote the resulting off-diagonal kernel by $K_f$.
Subtract the equilibrium value of each coefficient one factor at a time.
Terms outside the corner layer give $C\kappa$ times the weighted absolute
kernel. If the bad factor is a fixed input or output, the resulting
operator contains $\chi_{E_\eta}\widetilde K_{i,N}$ or
$\widetilde K_{i,N}\chi_{E_\eta}$, whose Hilbert--Schmidt norm tends to zero.
If the bad factor is on an unfixed leg, retain its indicator inside the
spherical or planar fiber. For almost every pair of fixed legs, the
fiber area of this indicator tends to zero; the exceptions are only
degenerate diagonals or coordinate lines.
The weighted Hilbert--Schmidt kernel provides a majorant, so another
application of dominated convergence gives
\[
 \|\nu_N^{-1/2}(K_f-K_N)\nu_N^{-1/2}\|_{2\to2}
                  \le C\kappa+o_{\eta\to0}(1).
\]
All errors are uniform on compact parameter sets.

It remains to exclude escape of mass into the corners.
If the conclusion failed, choose $\kappa_j,\eta_j\to0$ and unit vectors
$v_j=\sqrt{\nu_{N_j}}h_j$, with $P_{N_j}h_j=0$, for which the
left-hand side tends to zero or is negative.
Pass to a subsequence with $N_j\to N$ and $v_j\rightharpoonup v$.
The dual null-space vectors
$\kappa_{N_j}/\sqrt{\nu_{N_j}}$ converge strongly in $L^2$, so $v$
still satisfies the weighted microscopic constraint.
Compactness gives convergence of the kernel quadratic forms.
Set $\alpha_j=\nu_{f_j}/\nu_{N_j}$.
Since $\alpha_j\ge a_*>0$, the sequence is uniformly bounded, and
$\alpha_j\to1$ almost everywhere,
\[
 \liminf\int\alpha_j|v_j|^2
       \ge\|v\|_2^2+a_*(1-\|v\|_2^2).
\]
This follows by expanding $v_j=v+(v_j-v)$ and applying dominated
convergence and weak convergence.
Proposition \ref{geom:coercivity} therefore bounds the limiting energy
below by $c\|v\|_2^2+a_*(1-\|v\|_2^2)>0$, a contradiction.
\end{proof}

Finally, consider the products arising when the collision operator is
differentiated in space. Apply spatial interpolation to each fixed momentum
quadruple before integrating over the resonance measure.
The same measure retains the correlations between the legs; the required
momentum estimates reduce to its single-leg marginals and two-leg
conditional averages.

\begin{lemma}[Third-order products on the same four legs]\label{geom:moser}
Set $\dd\mu_* =\prod_i n_*(k_i)\dd\Xi_D$ and
$\dd m_*(k)=n_*^2\nu_*\dd k$. Every leg has marginal $m_*$.
For $i\ne j$, the conditional averaging operator $P_{ij}$ is determined by
\[
 \int_D g(k)(P_{ij}h)(k)\dd m_*(k)
    =\int g(k_i)h(k_j)\dd\mu_*
\]
for all bounded $g,h$.
It is a positive contraction on each $L^p(m_*)$; set
\[
 \omega_\eta=\max_{i\ne j}
       \|\chi_{E_\eta}P_{ij}\chi_{E_\eta}\|_{2\to2}\longrightarrow0.
\]
Suppose that a real function $e$ on $\T^3_X\times D$ satisfies
$|e|\le K_0$ everywhere and $|e|\le\kappa\le1$ outside the layer.
For $l=2,3$, if the total spatial derivative order is $l$ and
at least two distinct $e$ legs carry derivatives of positive order,
then every product of two or three factors satisfies
\begin{equation}\label{geom:localizedmoser}
 \left\|\prod_j\partial_X^{\beta_j}e(X,k_{i_j})\right\|_{L^2_XL^2(\mu_*)}
 \le C_{K_0}(\kappa^{1/3}+\omega_\eta^{1/3})
                          \|e\|_{H^l_XH_\nu}.
\end{equation}
No statistical independence of the legs is assumed.
General products of total order $l\le3$ satisfy the bound
$C_{K_0}\|e\|_{H^l_XH_\nu}$ without a small coefficient.
If the total order is at most three and at least one factor is a
derivative of a parameter error $a(X)$, then, in the range
$\|a\|_{H^3_X}\le A_0$, the corresponding product of scalar and
Hilbert-valued functions is bounded by
$C_{A_0}\|a\|_{H^3_X}\|e\|_{H^3_XH_\nu}$.
The remaining coefficients are bounded derivatives of fixed finite
order of prescribed functions, or bounded zeroth-order factors $e$.
\end{lemma}
\begin{proof}
Equality of the four marginals gives the positive contraction property.
Its unitary representation on $L^2(m_*)$ is exactly
$\nu_*^{-1/2}K_i\nu_*^{-1/2}$, so Lemma \ref{geom:compact}
implies that the corner-restricted norm tends to zero.
Interpolation with the $L^1,L^\infty$ contractions gives
\[
 \|\chi_{E_\eta} P_{ij}\chi_{E_\eta}\|_{p\to p}
       \le\omega_\eta^{2\min(1/p,1-1/p)}.
\]
For a periodic function at fixed momentum, integration by parts and
Hölder's inequality give
\[
 \|\nabla e\|_4^2\le C\|e\|_\infty\|D^2e\|_2,
 \quad \|\nabla e\|_6^2\le C\|e\|_\infty\|D^2e\|_3,
 \quad \|D^2e\|_3^2\le C\|\nabla e\|_6\|D^3e\|_2.
\]
Eliminating between the last two inequalities yields
$\|\partial^je\|_{6/j}\le C\|e\|_\infty^{1-j/3}
\|e\|_{H^3}^{j/3}$ for $j=1,2$; the second-order estimate is similar.
At zeros of the Hessian norm, one may first regularize the square root
and then pass to the limit.

Apply these estimates in $X$ for a fixed momentum quadruple, and then
apply Hölder's inequality on the same measure $\mu_*$.
For the third-order distribution $2+1$, the squared integral leaves
$F_i^{2/3}F_j^{1/3}$, where $F_i=\|e(\cdot,k_i)\|_{H^3}^2$.
If a differentiated leg lies outside the layer, its small amplitude
supplies at least $\kappa^{1/3}$.
If both legs lie inside the layer, the conditional-average bound at
$p=3$ or $3/2$ supplies $\omega_\eta^{2/3}$ in the squared integral.
For the third-order distribution $1+1+1$, write
$\chi_i=\mathbf1_{E_\eta}(k_i)$ and use
\[
 \int\chi_i\chi_j\chi_k F_i^{1/3}F_j^{1/3}F_k^{1/3}\dd\mu_*
 \le\left(\int\chi_i\chi_j\sqrt{F_iF_j}\dd\mu_*\right)^{2/3}
     \left(\int\chi_kF_k\dd\mu_*\right)^{1/3},
\]
followed by the $L^2$ corner-restricted bound for the two-leg conditional
average. Taking square roots gives the asserted exponent $1/3$ in
both cases. The second-order distribution $1+1$ gives the stronger
exponent $1/2$; zeroth-order legs are bounded by $K_0$.
Without corner restrictions, the same calculation using only contraction
gives the general Moser bound.

Finally, abbreviate spatial derivatives of total order $j$ by $a_j,e_j$;
these subscripts denote derivative order, not collision legs.
The highest-order distributions are $a_3e_0,a_2e_1,a_1e_2$.
Use, respectively,
$L^2_X\times L^\infty_XH_\nu$,
$L^6_X\times L^3_XH_\nu$, and
$L^\infty_X\times L^2_XH_\nu$.
If an additional $e$ leg occurs, first use its zeroth-order bound or the
same lower-order four-leg Moser estimate; for example, $a_1e_1e_1$
is controlled by $a_1\in L^\infty$ and the second-order product bound.
Multiple parameter derivatives are treated by the same distributions
of total order at most three.
The entire proof uses only $H^3_XH_\nu$ and the joint zeroth-order
amplitude bound, not $\sup_kH^3_X$ or $H^3_XL^\infty_k$.
\end{proof}
